\label{sub:IncreasingTheSizeOfExistingDPFs}
Another rather trivial approach is hardcoding larger datatypes replacing the original datatypes.
The user chooses her preferred data-size, specifies it before compilation, and uses this specified size during exuction.
This will commit each duoram to a single size which, however, can be chosen freely before compilation.
The code supports 64, 128, 256, and 532 bit values. %TODO das stimmt ja gar nicht
The implementation is located on branch \code{64ToUnlimited} in the author's repository.

\subsubsection*{Locating necessary Changes} \label{subsub:necessaryChanges}
Considering the source code's structure, the author chose a top-down approach for implementing this solution.
This includes locating the most general datatypes, distinguishing cases and assigning each case its special datatype.
The most optimal solution could trigger a trickle-down-effect, where more general datatypes affect more specialized datatypes with the ultimate goal of changing a single datatype or even a single flag in order to adjust the whole process of storing and retreiving data of a special size.
Therefore, the required efforts of achieving the optimal solution for this approach extend the efforts of solely switching the datatype used by the RDPFs.

The topmost elements in  this datatype-top-to-bottom structure are macros (\ref{subsub:implementationTypeshpp}).
These are specified by \code{types.hpp}.
They specify the most general datatypes and thus their derivates.
The most important alias, \code{value\_t}, by itself is used 269 times throughout the code.
Its most prominent usage is in the register types \code{RegAS} and \code{RegXS}, and the size of DPF-leafs.

As already explained in \ref{subsub:implementationTypeshpp}, \code{types.hpp} also holds \code{ADDRESS\_MAX\_BITS} specifying the size of addresses impacting the used datatype.
The derived type-alias \code{address\_t} is used 143 times throughout the base-implementation.

\subsubsection*{Macro Case-Distinctions}\label{subsub:macroCaseDistinction}
Cpp offers functionality for case distinction via macros.
This enables to use or not to use code based on the set macro.


It is set using the \code{VALUE\_BITS} macro, implementing a case distinction based on the value set for \code{VALUE\_BITS}.

\subsubsection*{Arbitrary large AES}\label{subsub:arbitraryLargeAES}
Each label in a DPF tree is AES encrypted.
This affects generation and traversion of such trees.
The base implementation uses an approach based on Intel AES-NI (\ref{subsub:implementationAesni}).
It is solely implemented for \code{\_\_m128i} vectors.
This fits the base definition of 128-bit labels in a DPF-tree.

Scaling the duoram requires the expansion to larger datatypes.
Current Intel AES-NI implementations, however, are solely described for 128, and 256 bit values. %TODO hier vielleicht das tolle paper verlinken und stimmt das überhaupt?
This contradicts the goal of scaling duoram arbitrarily.
Therefore, the author chose a solution similar to the prior approach of combination while leaving AES' security premises untouched.
By separating the base input in 128-bit blocks, encrypting each block independtly, and concatenating them after encryption makes AES scalable.
Thus, by looping the encryption dependent on the size of the given input, one implements an arbitrarily scalable AES protocol.

\subsubsection*{Adjusting Compilation}\label{subsub:adjustingCompilation}
To switch the datatype before compilation one needs to set a flag bit.
\code{-DVALUE\_BITS} is set in the \code{MAKEFILE}.
It represents the length of data one wants to store and retreive.
It will trigger the \code{VALUE\_BIT} macro and place the given value into it.
\code{VALUE\_BIT} then automatically adjusts the other macros which then will change the used datatypes and protocols to match the given value.
If nothing else is specified, \code{-DVALUE\_BITS} is set to $128$ by default which makes the code use \code{\_\_m128i} vectors for $M$.

\subsubsection*{Difficulties}

\subsubsection*{General Discussion}